\documentclass[11pt]{article}
\usepackage[margin=1in]{geometry}
\usepackage{amsmath,amssymb,amsfonts,mathtools}
\usepackage{array,booktabs,enumitem,graphicx,xcolor,hyperref}
\usepackage[T1]{fontenc}
\usepackage[utf8]{inputenc}
\title{On the Perturbations of String-Theoretic Black Holes}
\author{Gerald Gilbert}
\date{}
\begin{document}
\maketitle
\begin{abstract}

The perturbations of string-theoretic black holes are analyzed by generalizing the method of Chandrasekhar. Attention is focussed on the case of the recently considered charged string-theoretic black hole solutions as a representative example. It is shown that string-intrinsic effects greatly alter the perturbed motions of the string-theoretic black holes as compared to the perturbed motions of black hole solutions of the field equations of general relativity, the consequences of which bear on the questions of the scattering behavior and the stability of string-theoretic black holes. The explicit forms of the axial potential barriers surrounding the string-theoretic black hole are derived. It is demonstrated that one of these, for sufficiently negative values of the asymptotic value of the dilaton field, will inevitably become negative in turn, in marked contrast to the potentials surrounding the static black holes of general relativity. Such potentials may in principle be used in some cases to obtain approximate constraints on the value of the string coupling constant.

\end{abstract}

\section{Introduction}

='177 ='177 ='177 ='177 ='177 ='177 ='177 ='177 ='177 ='177 ='177 \textbackslash{}@magscale\#1 scaled \#1 = ammi10 \textbackslash{}@magscale5 ='177 ='60 ='60 ='60 ='60 ='60 ='60 ='60 ='60 ='60 ='60 ='60 = amsy10 \textbackslash{}@magscale5 ='60 `@=11 \#1 \textbackslash{}@@underline\#1 math expression `@=12 = = = = = = = = = =

\section{Method}

=cmex10 scaled 1200 \#1 \textasciicircum{} \#1 \#1 \_ \#1 \#1 /\#1 \#1 \#1 \#1 \#1 \#1 \#1 \#1 \#1 \#1 \#1 \#1 \#1 \#1 \#1 \#1 | \#1 | \#1 \#1 \#1 \#1 | \#1 |

\begin{equation}
\label{eq:compact}
L(\theta) = \sum_{i=1}^{n} \ell(x_i, \theta)
\end{equation}
Equation~\ref{eq:compact} is included to keep the sample paper-like while remaining small.

\section{Conclusion}

\#1\#2 1 .5ex 1 .4ex \#1\#2 1 .5ex 1 .3ex \#1\#2 1 .5ex \#2 \#1\#2 \#1 \#2 \#1\#2 \#1 \#2 \#1\#2\#3 \textasciicircum{}2 \#1 \#2 \#3 \#1\#2 \textasciicircum{} \#1 \#1 \#1 =0pt plus 1000pt minus 1000pt
\end{document}
